\section{A boundary deformation in three dimensions}\label{n3:intro:section}

The aim of this part is to replace an initial-data exterior by a
Riemannian exterior while keeping a genuine inner boundary. Three
quantities have to be controlled together: scalar curvature, the area
of every enclosing cut, and the ADM energy. The inner boundary creates
a flux term in the energy balance. We choose its geometry and the
boundary conditions so that this term is absorbed by a positive bulk
integral.

\subsection{Initial data, normals, and enclosing area}

We work in three spatial dimensions, with zero cosmological constant and
\(G=c=1\). For a Riemannian metric \(g\) and a symmetric covariant
two-tensor \(K\), define the constraint densities by
\begin{equation}\label{n3:intro:constraints}
 16\pi\mu=R_g+(\tr_gK)^2-|K|_g^2,\qquad
 8\pi J=\Div_g\bigl(K-(\tr_gK)g\bigr).
\end{equation}
Here \(J\) is a covector field. The dominant energy condition is
\begin{equation}\label{n3:intro:dec}
                         \mu\ge |J|_g.
\end{equation}
All initial data considered here are smooth, including at a compact
boundary when one is present.

On an asymptotically flat end with Euclidean coordinates \(x\) and
\(r=|x|\), our basic assumptions are
\begin{equation}\label{n3:intro:decay}
 g_{ij}-\delta_{ij}=O_2(r^{-q}),\qquad
 K_{ij}=O_1(r^{-1-q}),\qquad q>\tfrac12.
\end{equation}
The notation \(O_j(r^{-a})\) includes the corresponding coordinate
derivative bounds through order \(j\). We assume that \(\mu\) and
\(|J|_g\) are integrable and that the following ADM limits exist and are
finite:
\begin{align}
 E&=\frac1{16\pi}\lim_{r\to\infty}
   \int_{S_r}(\partial_jg_{ij}-\partial_i g_{jj})n^i
                    \dd A_\delta,\label{n3:intro:energy}\\
 P_i&=\frac1{8\pi}\lim_{r\to\infty}
   \int_{S_r}\bigl(K_{ij}-(\tr_gK)g_{ij}\bigr)n^j
                    \dd A_\delta.\label{n3:intro:momentum}
\end{align}
The normal and area form in these definitions are Euclidean. Whenever
\(E>|P|_\delta\), write
\begin{equation}\label{n3:intro:rest-mass}
                         m=\sqrt{E^2-|P|_\delta^2}.
\end{equation}

For a two-sided surface \(\Sigma\) with specified unit normal \(\nu\),
we use
\[
 H_\Sigma=\Div_\Sigma\nu,\qquad
 \theta_+(\Sigma)=H_\Sigma+\tr_\Sigma K.
\]
Thus Euclidean spheres have positive mean curvature for the normal
toward infinity. A future marginally outer trapped surface, abbreviated
MOTS, satisfies \(\theta_+=0\). A weakly future outer trapped surface
satisfies \(\theta_+\le0\). Our spacetime sign convention is
\begin{equation}\label{n3:intro:k-sign}
 K(Y,Z)=\mathbf g(\mathbf\nabla_Y n,Z),
\end{equation}
where \(n\) is the future unit timelike normal. Equivalently, in Gaussian
normal coordinates the spatial metric has normal derivative \(2K\).

\begin{definition}[Exterior and enclosing area]\label{n3:intro:exterior-class}
An exterior is a connected orientable smooth three-manifold \(\Omega\)
with nonempty compact smooth boundary, complete as a metric space with
its boundary included, and with one asymptotically flat coordinate end
whose complement is compact.
An enclosing cut is the entire compact smooth embedded two-sided
intrinsic boundary of a connected smooth end-containing domain $D$
closed in $\overline\Omega$, whose intrinsic interior lies in
$\operatorname{int}\Omega$. It separates the entire inner boundary
from the distant end. The inner side includes the original obstacle
and every bounded complementary pocket. All components of the cut
and every portion coincident with the inner boundary are counted.
Its normal points into $D$, toward the end. We put
\[
 a_g(\partial\Omega)=
  \inf\{|\Gamma|_g:\Gamma\text{ is an enclosing cut}\}.
\]
When using perimeter compactness, we take the closure of this class
among filled inner sets and include the area of any frontier coinciding
with the obstacle. Smooth outward approximation gives the same
infimum. Enclosing cuts and minimizing sets are taken with bounded
complementary pockets filled. Auxiliary perimeter arguments may use
competitors containing the inner obstacle before filling: filling a bounded
pocket deletes its frontier, cannot increase perimeter, and leaves the enclosing
infimum unchanged.
\end{definition}

The definition permits coincidence because first variations of the
minimum need not be realized by a cut lying strictly outside the
obstacle. It permits disconnected cuts because neither the minimizing
hull nor an intermediate trapped boundary need be connected. The smooth-cut and full-perimeter formulations have the same infimum;
we use the independent full-enclosure geometry of~\cite[Lemma~2.1]{RiemannianCompanion}.

\subsection{Prepared data and the geometric output}

The construction itself starts from the following explicit geometric
conditions. Its hypotheses do not require that the data have been
obtained by a particular end-replacement procedure.

\begin{definition}[Prepared three-dimensional data]\label{n3:prepared:data}
Let $(\Omega_0,g,K)$ be an exterior in
Definition~\ref{n3:intro:exterior-class}, with inner boundary $S_0$.
Suppose $K$ is compactly supported and, on its single coordinate end,
$g-\delta=O_j(r^{-1})$ for every fixed $j\ge0$. Assume the physical
constraint densities in Equation~\eqref{n3:intro:constraints} are
integrable, the ADM energy is $E>0$, and the ADM momentum is zero.
There are $0<\delta<1$, $c>0$, and a smooth radius extension
$\varrho\ge1$, equal to $r$ sufficiently far out, such that
\[
 R_g=O(r^{-3-\delta}),\qquad
 8\pi(\mu-|J|_g)\ge c\varrho^{-3-\delta}.
\]
The boundary is strictly future trapped with its normal into the
exterior: $\theta_+(S_0)<0$ on every component. Its full enclosing-area
infimum is denoted by $A_*=a_g(S_0)>0$.
\end{definition}

\begin{theorem}[Three-dimensional boundary deformation]
\label{n3:prepared:output}
For every prepared exterior $(\Omega_0,g,K)$ there is a connected
one-ended subexterior $\Omega\subset\Omega_0$ whose smooth compact
boundary $B$ separates $S_0$ from infinity, with the following property.
For each sufficiently small fixed $\epsilon>0$ and every sufficiently
large integer $N$, there is a smooth metric $\widehat g_{\epsilon,N}$
on $\overline\Omega$, complete with its boundary included, for which
\begin{align}
 R_{\widehat g_{\epsilon,N}}&\ge0,
 &H_{\partial\Omega}(\widehat g_{\epsilon,N})&<0,
 \label{n3:prepared:curvature}\\
 \widehat g_{\epsilon,N}&\ge e^{-4\epsilon}g,
 &E_{\widehat g_{\epsilon,N}}&\le E+\eta_{\epsilon,N},
 \label{n3:prepared:comparison}
\end{align}
where the boundary normal points into $\Omega$ and, at fixed
$\epsilon$, $\eta_{\epsilon,N}\to0$. More precisely, for constants
$C,a$ depending on the fixed data and $\epsilon$,
\[
 0\le\eta_{\epsilon,N}\le
 C(1+N)^a\bigl(e^{-\epsilon N}+e^{-\epsilon N/2}\bigr).
\]
The new end has $\widehat g_{\epsilon,N}-\delta=O_2(r^{-1})$ and
$R_{\widehat g_{\epsilon,N}}=O(r^{-3-\delta'})$ for some
$\delta'>0$. Every full enclosing cut $T$ in $\Omega$ satisfies
\[
 |T|_{\widehat g_{\epsilon,N}}\ge e^{-4\epsilon}A_*.
\]
\end{theorem}

The construction begins by selecting maximal black and white trapped
regions. Their outer collars prescribe opposite signs for the graph
slope, but the same negative mean-curvature condition for the final
metric. Section~\ref{n3:domains:section} proves the required geometry,
including the passage from weak supports to smooth enclosures.
Section~\ref{n3:deformation:section} specifies the trace and divergence
equations and excludes a lower conformal barrier. The estimates of
Section~\ref{n3:apriori:section} then control the solutions before
uniform ellipticity is available. Section~\ref{n3:existence:section}
supplies the boundary solve and degree argument. Finally,
Section~\ref{n3:limit:section} computes the end charge and proves
Theorem~\ref{n3:prepared:output}.

The algebraic inputs from~\cite[Propositions~4.1--4.2 and Lemma~4.3]{NumericalCompanion} are its scalar
identity, weighted coercivity, and the pointwise algebraic clause of
its floor lemma. Their
proofs are independent of its numerical conclusion. The boundary
construction here retains the signed-collar estimates, boundary flux,
homotopy, and degree that a filled construction does not require.

\subsection{Consequence for weakly decaying original data}

\begin{theorem}[Minimum-enclosing-area inequality]\label{n3:intro:exterior-inequality}
Let \((\Omega,g,K)\) be an exterior as in
Definition~\ref{n3:intro:exterior-class}. Assume
Equations~\eqref{n3:intro:constraints}--\eqref{n3:intro:decay}, integrability
of \(\mu\) and \(|J|_g\), and the finite ADM limits
\eqref{n3:intro:energy}--\eqref{n3:intro:momentum}. Orient \(\partial\Omega\)
by the normal into \(\Omega\), and suppose
\(\theta_+(\partial\Omega)\le0\). If \(E>|P|_\delta\), then
\begin{equation}\label{n3:intro:penrose}
          \sqrt{E^2-|P|_\delta^2}
               \ge \sqrt{\frac{a_g(\partial\Omega)}{16\pi}}.
\end{equation}
No extension of the data across \(\partial\Omega\) is required.
\end{theorem}


The last assertion will follow from the boundary construction and the
weak three-dimensional preparation theorem of~\cite[Proposition~7.13]{NumericalCompanion}.
That preparation supplies compact tensor support, strict trapping and
the weighted margin in Definition~\ref{n3:prepared:data}; its energies
converge to $\sqrt{E^2-|P|^2}$ and its full enclosing areas converge
to the original infimum. We will use only this preparation theorem,
not the companion's numerical inequality. In particular, the ADM
normalizations in Equations~\eqref{n3:intro:energy} and
\eqref{n3:intro:momentum} remain unchanged.
